\section{Details on ICU Mortality Example}
\label{appendix:icu_details}

This appendix provides full details on how we constructed the ICU subgroups, prepared the analytic dataset, estimated the OBD, identified sign reversals, and computed bootstrap standard errors for all subgroups used in Section~\ref{sec:real_example}. The analyses use the PhysioNet ICU cohort collected for the study of in-hospital mortality \citep{goldberger2000physiobank}. The raw cohort contains routine admission measurements for $3,997$ patients with known gender ($2,246$ women and $1,751$ men). We merge these records with the corresponding mortality outcomes and restrict attention to patients with complete information on the covariates used in the decomposition: age, ICU type, heart rate (HR), mean arterial pressure (MAP), temperature, urine output, and the indicator for in-hospital death. After removing observations with missing covariates, the final analytic dataset contains $3,429$ patients. All subgroup definitions, model specifications, and statistical procedures described below are applied to this cleaned dataset.

\subsection*{Model specification}

For each subgroup, we estimate separate linear models for men and for women.  
Let $g \in \{\text{men}, \text{women}\}$.  
We model the conditional expectation of in-hospital mortality as
\[
    \E[Y \mid X, g] \;=\; \alpha_{g} + X^{\top}\beta_{g},
\]
where $Y$ is the indicator for in-hospital mortality;  
$X$ is the vector of admission covariates (heart rate, mean arterial pressure, temperature, urine output, age, and ICU type);  
$\alpha_{g}$ is the group-specific intercept;  
and $\beta_{g}$ is the group-specific coefficient vector.  
The parameters $(\alpha_g,\beta_g)$ are estimated using ordinary least squares.  
Although $Y$ is binary, this follows the standard Oaxaca--Blinder framework, which relies on linear prediction rules.

\subsection*{Initial analysis}

Our first step was to evaluate whether sign reversals appeared in coarse, clinically meaningful partitions of the data. Specifically, we evaluated the OBD across nine predefined subgroups: the full cohort, each ICU type category, two broad age groups (patients younger than $65$ and patients aged $65$ or older), and two groups defined by a median split of the SAPS--I score. For each of these subgroups, we estimated the gender specific linear models described above and computed the explained and unexplained components under both reference rules. As reported in \cref{tab:A1_initial_subgroups}, none of these nine subgroups exhibited a sign reversal in the explained or unexplained components. This motivated a more aggressive and systematic search for sign flips, which we describe next.
\begin{table}[H]
    \centering
    \small
    \caption{OBD for the nine predefined ICU subgroups}
    \label{tab:A1_initial_subgroups}
    \begin{tabular}{l|r|r|cc|cc}
    \hline
    \textbf{subset} & \textbf{n} & \textbf{Total gap} 
        & \multicolumn{2}{c|}{\textbf{Men reference}} 
        & \multicolumn{2}{c}{\textbf{Women reference}} \\ 
    \cline{4-7}
      &   &   
      & \textbf{Explained} & \textbf{Unexplained} 
      & \textbf{Explained} & \textbf{Unexplained} \\ \hline
    All patients    & 3429 &  0.0079 &  0.0114 & -0.0035 &  0.0156 & -0.0077 \\ 
    ICUType 1.0     &  490 &  0.0295 &  0.0306 & -0.0012 &  0.0368 & -0.0073 \\ 
    ICUType 2.0     &  643 & -0.0010 &  0.0039 & -0.0049 &  0.0040 & -0.0050 \\ 
    ICUType 3.0     & 1357 & -0.0218 &  0.0043 & -0.0261 &  0.0139 & -0.0357 \\ 
    ICUType 4.0     &  939 &  0.0202 &  0.0248 & -0.0046 &  0.0208 & -0.0005 \\ 
    Age $< 65$      & 1567 &  0.0170 &  0.0085 &  0.0085 &  0.0055 &  0.0115 \\ 
    Age $\ge 65$    & 1862 & -0.0098 &  0.0101 & -0.0199 &  0.0159 & -0.0256 \\ 
    SAPS $\le 15$   & 1971 &  0.0257 &  0.0093 &  0.0164 &  0.0121 &  0.0136 \\ 
    SAPS $> 15$     & 1458 & -0.0157 &  0.0104 & -0.0261 &  0.0133 & -0.0290 \\ \hline
    \end{tabular}
\end{table}


\subsection*{Subgroup construction and identification of the HR sign flip}

After verifying that none of the nine predefined ICU subgroups exhibited a sign reversal in the explained or unexplained OB components, we expanded the analysis to a much larger and more heterogeneous collection of subgroups. The goal of this second stage was not to target any particular physiological pattern, but rather to explore systematically whether sign reversals arise in realistic data partitions that a practitioner might naturally inspect.

We constructed $139$ subgroups in total. These sets include: quartiles of each continuous admission vital (heart rate, MAP, temperature, urine output);  
deciles of age; ICU-type specific cohorts; clinically motivated threshold subsets (HR $>100$, MAP $<65$, Temp $> 38^\circ$C, Urine $>1000$ mL); fifty random $50\%$ subsamples of the cohort; fifty random $30\%$ subsamples of the cohort.

Across all $139$ subgroups, we found that $27$ exhibited at least one sign reversal. Many of these arose within the random subsamples, indicating that the OBD can yield unstable conclusions even when applied to simple resampled partitions of the same population. Several physiologically interpretable groups also displayed sign changes. The HR quartile two subgroup examined in Section~\ref{sec:real_example} is one such case, and serves as a concrete illustration of how these reversals can appear in clinically meaningful settings.

\textbf{Bootstrap}
For each of the $27$ subgroups that exhibited a sign reversal, we computed bootstrap standard errors to assess the stability of the OB components. The bootstrap procedure follows the standard Oaxaca--Blinder framework: men and women are resampled \emph{within} group with replacement, preserving group sizes. For each bootstrap replication $b=1,\dots,B$, we estimate the OB decomposition using the same linear models described above.


Let $\hat\theta$ denote any Oaxaca--Blinder component (for example, the explained component under the women reference model), and let  
$\hat\theta^{(1)},\dots,\hat\theta^{(B)}$ denote its bootstrap replicates.  
The bootstrap standard error and mean are
\[
    \widehat{\operatorname{se}}(\hat\theta)
    =
    \sqrt{
    \frac{1}{B-1}\sum_{b=1}^B
    \bigl(\hat\theta^{(b)} - \bar\theta\bigr)^2
    },
    \qquad
    \bar\theta = \frac{1}{B}\sum_{b=1}^B \hat\theta^{(b)}.
\]

P-values were computed using a normal approximation,
\[
    p = 2\left(1 - \Phi\left(\frac{|\hat\theta|}{\widehat{\operatorname{se}}(\hat\theta)}\right)\right).
\]

For each subgroup with a sign flip, we report:  
(i) the estimated mortality gap (men minus women);  
(ii) the explained and unexplained components under the men reference rule;  
(iii) the explained and unexplained components under the women reference rule; and  
(iv) their corresponding bootstrap standard errors and significance levels.  
These complete results appear in Table~\ref{tab:icu_mortality}. The persistence of the sign reversals after accounting for sampling variability indicates that the instability documented in Section~\ref{sec:real_example} is not a numerical artifact, but an inherent behavior of the OBD in this clinical dataset.


\begin{table}[H]
    \centering
    \caption{ICU Mortality Gap Decomposition for Subgroups Exhibiting Sign Reversals}
    \label{tab:icu_mortality}
    \scriptsize
    \begin{tabular}{@{}lc*{5}{c}@{}}
    \toprule
    \multirow{2}{*}{\textbf{Subset}} &
    \multirow{2}{*}{\textbf{n}} &
    \textbf{Gap} &
    \multicolumn{2}{c}{\textbf{Female Reference}} &
    \multicolumn{2}{c}{\textbf{Male Reference}} \\
    \cmidrule(lr){3-3} \cmidrule(lr){4-5} \cmidrule(lr){6-7}
    & & (M–F) & Expl. & Unexp. & Expl. & Unexp. \\
    \midrule
    
    Age decile 1 & 344 & 0.027 & 0.002 & 0.025 & 0.036$^{*}$ & $-0.009$ \\
    & & (0.033) & (0.009) & (0.033) & (0.019) & (0.030) \\[0.3em]
    
    Age decile 2 & 347 & 0.004 & 0.023 & $-0.019$ & $-0.010$ & 0.015 \\
    & & (0.033) & (0.021) & (0.041) & (0.014) & (0.036) \\[0.3em]
    
    Age decile 4 & 372 & $-0.015$ & 0.012 & $-0.027$ & $-0.004$ & $-0.012$ \\
    & & (0.032) & (0.013) & (0.035) & (0.017) & (0.037) \\[0.3em]
    
    Age decile 5 & 344 & $-0.025$ & 0.010 & $-0.035$ & $-0.008$ & $-0.017$ \\
    & & (0.039) & (0.016) & (0.039) & (0.022) & (0.042) \\[0.3em]
    
    Age decile 9 & 325 & $-0.077$ & 0.001 & $-0.078$ & 0.001 & $-0.078$ \\
    & & (0.048) & (0.023) & (0.049) & (0.013) & (0.049) \\[0.3em]
    
    \textbf{HR quartile 2} & 816 & 0.035 & 0.021$^{***}$ & 0.014 & $-0.007$ & 0.042$^{*}$ \\
    & & (0.023) & (0.007) & (0.025) & (0.011) & (0.025) \\[0.3em]
    
    MAP quartile 0 & 877 & $-0.014$ & $-0.003$ & $-0.011$ & 0.004 & $-0.018$ \\
    & & (0.026) & (0.009) & (0.026) & (0.008) & (0.026) \\[0.3em]
    
    MAP quartile 1 & 840 & 0.024 & 0.036$^{***}$ & $-0.012$ & 0.011 & 0.013 \\
    & & (0.022) & (0.012) & (0.027) & (0.008) & (0.022) \\[0.3em]
    
    Urine quartile 0 & 859 & $-0.031$ & 0.007 & $-0.038$ & $-0.006$ & $-0.025$ \\
    & & (0.027) & (0.009) & (0.027) & (0.007) & (0.027) \\[0.3em]
    
    MAP $< 65$ & 697 & $-0.028$ & $-0.011$ & $-0.017$ & 0.002 & $-0.030$ \\
    & & (0.029) & (0.012) & (0.029) & (0.008) & (0.029) \\[0.3em]
    
    Temp $> 36.5$ & 267 & $-0.043$ & $-0.001$ & $-0.042$ & 0.008 & $-0.051$ \\
    & & (0.047) & (0.025) & (0.050) & (0.019) & (0.046) \\[0.3em]
    
    Urine $> 1000$ & 210 & 0.046 & $-0.014$ & 0.060 & 0.014 & 0.031 \\
    & & (0.036) & (0.018) & (0.038) & (0.023) & (0.037) \\[0.3em]
    
    Random 50\% \#10 & 1714 & 0.012 & 0.013$^{***}$ & $-0.001$ & 0.011$^{***}$ & 0.000 \\
    & & (0.017) & (0.005) & (0.017) & (0.005) & (0.016) \\[0.3em]
    
    Random 50\% \#11 & 1714 & 0.012 & 0.014$^{***}$ & $-0.002$ & 0.009$^{***}$ & 0.003 \\
    & & (0.017) & (0.005) & (0.018) & (0.005) & (0.017) \\[0.3em]
    
    Random 50\% \#15 & 1714 & 0.009 & 0.010$^{***}$ & $-0.001$ & 0.007$^{***}$ & 0.001 \\
    & & (0.016) & (0.004) & (0.016) & (0.004) & (0.016) \\[0.3em]
    
    Random 50\% \#19 & 1714 & 0.013 & 0.014$^{***}$ & $-0.001$ & 0.009$^{***}$ & 0.004 \\
    & & (0.017) & (0.005) & (0.017) & (0.004) & (0.016) \\[0.3em]
    
    Random 50\% \#25 & 1714 & 0.018 & 0.019$^{***}$ & $-0.000$ & 0.016$^{***}$ & 0.003 \\
    & & (0.017) & (0.006) & (0.019) & (0.005) & (0.017) \\[0.3em]
    
    Random 50\% \#29 & 1714 & 0.015 & 0.016$^{***}$ & $-0.001$ & 0.014$^{***}$ & 0.002 \\
    & & (0.017) & (0.005) & (0.018) & (0.005) & (0.017) \\[0.3em]
    
    Random 50\% \#37 & 1714 & 0.016 & 0.017$^{***}$ & $-0.001$ & 0.016$^{***}$ & 0.000 \\
    & & (0.017) & (0.005) & (0.018) & (0.005) & (0.016) \\[0.3em]
    
    Random 30\% \#1 & 1028 & 0.021 & 0.016$^{***}$ & 0.005 & 0.022$^{***}$ & $-0.001$ \\
    & & (0.020) & (0.007) & (0.022) & (0.008) & (0.019) \\[0.3em]
    
    Random 30\% \#3 & 1028 & 0.020 & 0.022$^{***}$ & $-0.001$ & 0.017$^{***}$ & 0.004 \\
    & & (0.021) & (0.007) & (0.022) & (0.006) & (0.019) \\[0.3em]
    
    Random 30\% \#7 & 1028 & 0.017 & 0.030$^{***}$ & $-0.013$ & 0.016$^{***}$ & 0.000 \\
    & & (0.022) & (0.008) & (0.024) & (0.008) & (0.022) \\[0.3em]
    
    Random 30\% \#17 & 1028 & 0.013 & 0.012$^{**}$ & 0.001 & 0.016$^{***}$ & $-0.003$ \\
    & & (0.022) & (0.007) & (0.023) & (0.007) & (0.021) \\[0.3em]
    
    Random 30\% \#21 & 1028 & 0.014 & 0.014$^{***}$ & $-0.000$ & 0.012$^{**}$ & 0.002 \\
    & & (0.022) & (0.007) & (0.023) & (0.007) & (0.021) \\[0.3em]
    
    Random 30\% \#23 & 1028 & 0.005 & 0.013$^{**}$ & $-0.008$ & 0.002 & 0.003 \\
    & & (0.021) & (0.007) & (0.022) & (0.007) & (0.021) \\[0.3em]
    
    Random 30\% \#35 & 1028 & 0.010 & 0.012 & $-0.002$ & 0.002 & 0.008 \\
    & & (0.022) & (0.008) & (0.024) & (0.005) & (0.022) \\[0.3em]
    
    Random 30\% \#44 & 1028 & $-0.006$ & 0.024$^{***}$ & $-0.030$ & $-0.002$ & $-0.004$ \\
    & & (0.022) & (0.008) & (0.023) & (0.012) & (0.024) \\
    
    \bottomrule
    \end{tabular}
    \vspace{0.5em}
    
    \footnotesize
    \textit{Note:} Standard errors in parentheses computed via bootstrap resampling with $1{,}000$ replications.  
    P-values based on a normal approximation.  
    Significance levels: * $p<0.10$, ** $p<0.05$, *** $p<0.01$.  
    The decomposition separates the gender mortality gap (men minus women) into explained and unexplained components using the Oaxaca--Blinder method.  
    Female Reference uses the model for women as the reference outcome structure; Male Reference uses the model for men.
    \end{table}
    